% Part:sets-functions-relations
% Chapter: size-of-sets
% Section: comparing-sizes

\documentclass[../../../include/open-logic-section]{subfiles}

\begin{document}

\olfileid{sfr}{siz}{car}

\olsection{Verzamelingen van verschillende grootte en de stelling van Cantor}

\begin{explain}
We hebben precies vastgelegd wat het betekent dat twee verzamelingen
even groot zijn. We kunnen ook precies vastleggen wat het betekent dat de
ene verzameling kleiner is dan de andere. Voor onze definitie van ``kleiner
dan (of gelijkmachtig met)'' is, in plaats van !!a{bijection} tussen de
verzamelingen, !!a{injection} van de eerste verzameling naar de tweede nodig.
Als zo'n functie bestaat, is de grootte van de eerste verzameling kleiner
dan of gelijk aan die van de tweede. Intuïtief garandeert !!a{injection} van
de ene verzameling naar de andere dat het beeld van de functie minstens
evenveel !!{element}s bevat als het domein, omdat geen twee !!{element}s van
het domein op hetzelfde !!{element} van het beeld worden afgebeeld.
\end{explain}

\begin{defn}
$A$ is \emph{niet groter dan}~$B$, geschreven als $\cardle{A}{B}$, precies
dan wanneer er !!a{injection} $f \colon A \to B$ bestaat.
\end{defn}

Deze relatie is duidelijk reflexief en transitief, maar niet symmetrisch
(dit laten we als opgave). We kunnen ook een begrip invoeren waarmee we
uitdrukken dat de ene verzameling (strikt) kleiner is dan de andere. 

\begin{defn}
$A$ is \emph{kleiner dan}~$B$, geschreven als $\cardless{A}{B}$, precies dan
wanneer er !!a{injection}~$f\colon A \to B$ bestaat, maar geen
!!{bijection}~$g\colon A \to B$, oftewel $\cardle{A}{B}$ en
$\cardneq{A}{B}$.
\end{defn}

Ook deze relatie is duidelijk irreflexief en transitief. (Dit laten we als
opgave.) Met deze notatie kunnen we zeggen dat een verzameling $A$
!!{enumerable} is precies dan wanneer $\cardle{A}{\Nat}$, en dat $A$
!!{nonenumerable} is precies dan wanneer $\cardless{\Nat}{A}$. Hiermee
kunnen we
\oliflabeldef{sfr:siz:nen-alt:thm:nonenum-pownat}{%
\olref[sfr][siz][nen-alt]{thm:nonenum-pownat}
opnieuw formuleren als de vaststelling dat
$\cardless{\Nat}{\Pow{\Nat}}$}{%
\olref[sfr][siz][nen]{thm:nonenum-pownat}
opnieuw formuleren als de vaststelling dat
$\cardless{\PosInt}{\Pow{\PosInt}}$}. Sterker nog, \citet{Cantor1892}
bewees dat deze laatste uitspraak \emph{volkomen algemeen} geldt:

\begin{thm}[Cantor]\ollabel{thm:cantor}
$\cardless{A}{\Pow{A}}$, voor elke verzameling $A$.
\end{thm}

\begin{proof}
De afbeelding $f(x) = \{x\}$ is !!a{injection} $f \colon A \to \Pow{A}$,
want als $x \neq y$, dan geldt wegens extensionaliteit ook
$\{x\} \neq \{y\}$, en dus $f(x) \neq f(y)$. Bijgevolg geldt
$\cardle{A}{\Pow{A}}$. 

\begin{editorial}
We geven het uitvoerige bewijs als \olref[nen]{sec} aanwezig is, en anders
het kortere bewijs dat aansluit bij \olref[nen-alt]{sec}.
\end{editorial}

\oliflabeldef{sfr:siz:nen:sec}{% 
  We tonen nu aan dat er geen surjectieve functie~$g\colon A \to
  \Pow{A}$ kan bestaan, laat staan een bijectieve functie, en dus dat
  $\cardneq{A}{\Pow{A}}$. Stel namelijk dat $g\colon A \to \Pow{A}$.
  Omdat $g$ totaal is, wordt elke $x \in A$ afgebeeld op een deelverzameling
  $g(x) \subseteq A$. We kunnen aantonen dat $g$ niet surjectief kan zijn.
  Daartoe definiëren we een deelverzameling~$\overline{A} \subseteq A$ die
  volgens haar definitie niet tot het beeld van~$g$ kan behoren. Stel
  \[
  \overline{A} = \Setabs{x \in A}{x \notin g(x)}.
  \]
  Omdat $g(x)$ voor elke $x \in A$ gedefinieerd is, is $\overline{A}$
  duidelijk een welgedefinieerde deelverzameling van~$A$. Toch kan zij niet
  tot het beeld van~$g$ behoren. Neem een willekeurige $x \in A$; we zullen
  aantonen dat $\overline{A} \neq g(x)$. Als $x \in g(x)$, dan geldt
  $x \notin g(x)$ niet, en dus volgt uit de definitie van~$\overline{A}$
  dat $x \notin \overline{A}$. Als $x \notin g(x)$, dan wordt, omdat
  $x \in A$, voldaan aan de definiërende eigenschap van~$\overline{A}$,
  en dus geldt $x \in \overline{A}$. Uit beide gevallen volgt voor de
  verzameling~$\overline{A}$ en elke $x \in A$ dat $x \in g(x)$ precies dan wanneer
  $x \notin \overline{A}$, en dus $g(x) \neq \overline{A}$. Met andere
  woorden, $\overline{A}$ kan niet tot het beeld van~$g$ behoren, in
  tegenspraak met de aanname dat~$g$ surjectief is.}{Rest nog aan te tonen
  dat $\cardneq{A}{\Pow{A}}$. Neem voor een bewijs uit het ongerijmde aan dat
  $\cardeq{A}{\Pow{A}}$, oftewel dat er een !!{bijection}
  $g \colon A \to \Pow{A}$ bestaat. Beschouw nu:
  \[
    D = \Setabs{x \in A}{x \notin g(x)}
  \]
  Merk op dat $D \subseteq A$, zodat $D \in \Pow{A}$. Omdat $g$
  !!a{bijection} is, bestaat er een $y \in A$ waarvoor $g(y) = D$. Maar nu
  geldt:
  \[
    y \in g(y) \text{ precies dan wanneer } y \in D
      \text{ precies dan wanneer } y \notin g(y).
  \]
  Dit is een tegenspraak; dus $\cardneq{A}{\Pow{A}}$.}{}
\end{proof}

\begin{explain}
\oliflabeldef{sfr:siz:nen:thm:nonenum-pownat}{Het is verhelderend om het
  bewijs van \olref{thm:cantor} te vergelijken met dat van
  \olref[nen]{thm:nonenum-pownat}. Daar hebben we aangetoond dat men bij elke
  lijst $Z_1$, $Z_2$, \dots, van deelverzamelingen van~$\PosInt$ een
  verzameling~$\overline{Z}$ van getallen kan construeren die gegarandeerd
  niet op de lijst voorkomt. Zij komt gegarandeerd niet op de lijst voor,
  omdat voor elke $n \in \PosInt$ geldt dat $n \in Z_n$ precies dan wanneer
  $n \notin \overline{Z}$. Er is zo steeds een getal dat !!a{element} is van
  een van $Z_n$ en $\overline{Z}$, maar geen !!{element} van de andere. Hier
  volgen we hetzelfde idee, met dit verschil dat de indices~$n$ nu
  !!{element}s van~$A$ zijn in plaats van~$\PosInt$. De verzameling
  $\overline{A}$ is zo gedefinieerd dat zij van~$g(x)$ verschilt voor elke
  $x \in A$, omdat $x \in g(x)$ precies dan wanneer $x \notin
  \overline{A}$. Ook nu is er steeds !!a{element} van~$A$ dat !!a{element}
  is van een van $g(x)$ en $\overline{A}$, maar niet van de andere. En zoals
  $\overline{Z}$ daarom niet op de lijst $Z_1$, $Z_2$, \dots, kan voorkomen,
  zo kan $\overline{A}$ niet tot het beeld van~$g$ behoren.}{}

\oliflabeldef{sfr:siz:nen-alt:thm:nonenum-pownat}{Het is verhelderend om het
  bewijs van \olref{thm:cantor} te vergelijken met dat van
  \olref[nen-alt]{thm:nonenum-pownat}. Daar hebben we aangetoond dat we bij
  elke lijst $N_0$, $N_1$, $N_2$, \dots, van deelverzamelingen van~$\Nat$
  een verzameling~$D$ van getallen kunnen construeren die gegarandeerd niet
  op de lijst voorkomt. Zij komt gegarandeerd niet op de lijst voor omdat
  $n \in N_n$ precies dan wanneer $n \notin D$, voor elke $n \in \Nat$. Hier
  volgen we hetzelfde idee, met dit verschil dat de indices~$n$ nu
  !!{element}s van~$A$ zijn in plaats van~$\Nat$. De verzameling $D$ is zo
  gedefinieerd dat zij van~$g(x)$ verschilt voor elke $x \in A$, omdat
  $x \in g(x)$ precies dan wanneer $x \notin D$.}{}

Dit bewijs kan ook goed worden vergeleken met het bewijs van de paradox van
Russell, \olref[sfr][set][rus]{thm:russells-paradox}. De stelling van Cantor
vormde namelijk de inspiratie voor Russells eigen paradox.
\end{explain}

\begin{prob}
  Toon aan dat er geen !!a{injection} $g\colon \Pow{A} \to A$ kan bestaan,
  ongeacht de verzameling~$A$. Aanwijzing: stel dat
  $g\colon \Pow{A} \to A$ !!{injective} is. Beschouw
  $D = \Setabs{g(B)}{B \subseteq A \text{ en } g(B) \notin B}$. Stel
  $x = g(D)$. Leid een tegenspraak af uit het feit dat $g$ !!{injective} is. 
\end{prob}
\end{document}
